\documentclass[a4paper, 12pt]{article}

% --- PACCHETTI FONDAMENTALI ---
\usepackage[utf8]{inputenc} % Codifica caratteri (per le lettere accentate)
\usepackage[T1]{fontenc}    % Codifica font
\usepackage[italian]{babel} % Lingua italiana (sillabazione, capitoli, ecc.)

% --- PACCHETTI MATEMATICI (I "Must Have") ---
\usepackage{amsmath}  % Il cuore della matematica in LaTeX
\usepackage{amssymb}  % Simboli extra (R dei reali, C dei complessi...)
\usepackage{amsthm}   % Per teoremi, definizioni e dimostrazioni

% --- ESTETICA ---
\usepackage{geometry} % Per gestire i margini
\geometry{margin=2.5cm}
\usepackage{hyperref} % Per link cliccabili e indice interattivo

% --- DEFINIZIONE DI AMBIENTI PERSONALIZZATI ---
% Qui creiamo dei "box" per Teoremi e Definizioni
\newtheorem{teorema}{Teorema}
\theoremstyle{definition}
\newtheorem{definizione}{Definizione}
\newtheorem{esempio}{Esempio}

% --- INIZIO DEL DOCUMENTO ---
\title{\textbf{Appunti di Matematica e Analisi}}
\author{Studente Universitario}
\date{\today}

\begin{document}

\maketitle % Genera il titolo
\tableofcontents % Genera l'indice automatico
\newpage

\section{Algebra Lineare}

\subsection{Matrici e Autovalori}
Consideriamo una matrice $A \in \mathbb{R}^{2 \times 2}$. Il \textbf{polinomio caratteristico} è definito come:
\begin{equation} \label{eq:polinomio}
    p_A(\lambda) = \det(A - \lambda I)
\end{equation}
Le radici di questo polinomio sono gli \emph{autovalori} della matrice.

\begin{esempio}
Calcoliamo gli autovalori della matrice:
$$ A = \begin{pmatrix} 3 & 1 \\ 0 & 2 \end{pmatrix} $$
Applicando la formula \eqref{eq:polinomio}:
$$ \det \begin{pmatrix} 3-\lambda & 1 \\ 0 & 2-\lambda \end{pmatrix} = (3-\lambda)(2-\lambda) = 0 $$
Gli autovalori sono $\lambda_1 = 3$ e $\lambda_2 = 2$.
\end{esempio}

\section{Analisi Matematica}

\subsection{Limiti Notevoli}
Ricordiamo il limite fondamentale per il seno, valido per $x \to 0$:

$$ \lim_{x \to 0} \frac{\sin(x)}{x} = 1 $$

Da questo possiamo dedurre che per $\alpha > -1$:
$$ \lim_{x \to 0^+} x^\alpha \sin(x) = 0 $$

\subsection{Sviluppo di Taylor}
\begin{teorema}[Formula di Taylor con resto di Peano]
Sia $f: (a,b) \to \mathbb{R}$ una funzione derivabile $n$ volte in $x_0$. Allora:
$$ f(x) = \sum_{k=0}^n \frac{f^{(k)}(x_0)}{k!} (x-x_0)^k + o((x-x_0)^n) $$
\end{teorema}

Se applichiamo questo al coseno attorno a $x_0 = 0$:
\begin{align*}
    \cos(x) &= 1 - \frac{x^2}{2!} + \frac{x^4}{4!} + o(x^4) \\
            &= 1 - \frac{1}{2}x^2 + \frac{1}{24}x^4 + o(x^4)
\end{align*}
\textit{Nota: L'ambiente "align" permette di allineare le equazioni sul simbolo =.}

\section{Conclusioni}
Scrivere in \LaTeX\ permette di avere formule pulite come $\int